\section{Normalization and the finite-clock fibre}
\label{sec:polynomial}
The normalizing polynomial is constrained by the requirement that
$q(T_j)P(T_j)$ interpolate a matrix polynomial of degree at most $d$.
The operator in \eqref{eq:normalization} records precisely this constraint.

\begin{lemma}[The normalization defect]\label{lem:defect}
Under the rank assumption of Theorem~\ref{thm:main}, the greatest common
divisor of all entries of $K$ is $g=\gcd(f_1,\ldots,f_k)$. If
$\ell\ge2d+1$, then \eqref{eq:kernel} holds.
\end{lemma}
\begin{proof}
Every entry of $K$ is a constant linear combination of the $f_b$, so
$g$ divides those entries. Conversely, constant left and right channel
inverses recover $\diag(\alpha_b f_b)$. Every common divisor of the
entries therefore divides every $f_b$.

Let $r\in\ker\cN_P$ have degree less than $d$. There is a matrix
polynomial $H$ of degree at most $d$ with $H(T_j)=r(T_j)P(T_j)$.
Every entry of $qH-rK$ has degree at most $2d$ and vanishes at $2d+1$
nodes. Hence $qH=rK$ identically. Write $K=g\widetilde K$ and
$q=g\widetilde q$. The entries of $\widetilde K$ have greatest common
divisor one. Their Bezout identity implies $\widetilde q\mid r$.
Thus $r=\widetilde q s$ with $\deg s<\deg g$. Conversely each such
$s$ gives $H=\widetilde K s$ of degree at most $d-1$ and satisfies
the identity. This includes the case $\deg g=d$.
\end{proof}

Cancelling $g$ lowers the visible degree to $d-\deg g$. The zeros of
the cancelled factor cannot be read from the normalized curve. A shorter
design can have a normalization defect even when $g=1$.

\begin{lemma}[Positive local fibres]\label{lem:local-fibre}
Suppose $\ell\ge d+1$, the true channels have full column rank,
$\alpha_b>\alpha_*$, and all roots of each $f_b$ are simple in $(L,D)$.
If $0\ne r\in\ker\cN_P$, there is a smooth parameter curve with the
same $U,V$, unchanged observations at every sampled clock, and a
nonconstant root multiset.
\end{lemma}
\begin{proof}
Let $H$ interpolate $r(T_j)P(T_j)$ with degree at most $d$. Interpolation
at the first $d+1$ nodes shows that
$H=U\diag(h_b)V^{\mathsf T}$, with each $h_b$ of degree at most $d$.
Since every $P(T_j)$ has entry sum one, $\sum_bh_b-r$ vanishes at
$d+1$ nodes and is identically zero. For small real $u$ put
\[
 q_u=q+ur,\qquad K_u=K+uH,\qquad
 w_{b,u}=\alpha_bf_b+uh_b.
\]
Let $\alpha_{b,u}$ be the leading coefficient of $w_{b,u}$ and set
$f_{b,u}=w_{b,u}/\alpha_{b,u}$. Then $\sum_b\alpha_{b,u}=1$,
$q_u=\sum_b\alpha_{b,u}f_{b,u}$, and the strict weight inequalities
persist. Around each root of each component choose disjoint brackets
with opposite signs at their endpoints. Small coefficient perturbations
preserve a root in each bracket, accounting for all $d$ roots. These
real roots are smooth in $u$ by the implicit function theorem.
The identity $K_u(T_j)=q_u(T_j)P(T_j)$ proves exact agreement.

Suppose the aggregate multiset were constant on a neighbourhood of
$u=0$. Each continuously tracked component root would then take values
in a fixed finite set. A continuous map from a connected interval to
a finite set is constant. This argument applies also when roots of
different components coincide. Hence every $f_{b,u}=f_b$.
Applying the true channel inverses at a clock gives
$\alpha_{b,u}/\alpha_b=q_u(T_j)/q(T_j)$ for every $b,j$.
Summing the weights makes this ratio one. Thus $r$ vanishes at $d+1$
nodes, contradicting $r\ne0$ and $\deg r<d$.
\end{proof}

\begin{proposition}[Sharp worst-case clock count]\label{prop:clocks}
For every admissible fixed $k,m_1,m_2$, there are positive full-rank
channels and interior, relatively prime component polynomials for
which any specified $2d$ distinct clocks admit the fibre of
Lemma~\ref{lem:local-fibre}.
\end{proposition}
\begin{proof}
For $k=2$, write $E_b=U_{:b}V_{:b}^{\mathsf T}$ and
$P(T)=E_2+(E_1-E_2)w(T)$, where $w=\alpha_1f_1/q$.
Since $\Omega(r(T_j))_j=0$ for $\deg r<d$,
\[
 \cN_Pr=\big[\Omega(r(T_j)w(T_j))_j\big]
                   \otimes\operatorname{vec}(E_1-E_2).
\]
At $2d$ nodes the scalar residual space has dimension $d-1$, so the
map from the $d$ coefficients of $r$ has nonzero kernel, even for
relatively prime $f_1,f_2$.

For arbitrary $k$, choose close relatively prime monic polynomials
$f,g$ with simple roots in $(L,D)$, so that every convex combination
$(1-s)f+sg$ retains that property by the bracket argument. Choose
$s_b\in[0,1]$ including both endpoints and set
$f_b=(1-s_b)f+s_bg$, with $\alpha_b=1/k$ and any strictly positive
full-rank stochastic channels. The collection is relatively prime.
Put $A=\sum_b\alpha_bE_b$, $D_0=\sum_b\alpha_bs_bE_b$, and
$\bar s=\sum_b\alpha_bs_b$. Then
\[
 P(T)=A+(D_0-\bar sA)\frac{g(T)-f(T)}{q(T)}.
\]
This is a genuinely nonconstant affine-line curve. Indeed the matrices
$E_b$ are linearly independent, as follows by applying the channel
inverses. Thus $D_0-\bar sA\ne0$, since the $s_b$ are not all equal.
Moreover $g-f$ has degree less than $d$ and is nonzero, so $(g-f)/q$
is nonconstant. The same scalar residual-dimension argument gives a
nonzero kernel. Lemma~\ref{lem:local-fibre} completes the proof.
\end{proof}

\subsection{Root matching without cross-component gaps}
Write $\|F\|_{\rm coef}=\sum_{r=0}^d\|F_r\|_\op$ for a matrix
polynomial $F(z)=\sum_{r=0}^dF_rz^r$.
\begin{lemma}[Diagonal polynomial perturbations]\label{lem:roots}
Let $D_f(z)=\diag(\alpha_bf_b(z))$ be as in \eqref{eq:model}.
If $\|E\|_{\rm coef}\le\varepsilon$ is sufficiently small, then
$\det(D_f+E)$ has degree $kd$ and its root multiset $Z_E$ satisfies
$d_\infty(\Lambda,Z_E)\le C\varepsilon^{1/d}$.
The roots of $Z_E$ may be complex. If roots within each $f_b$ are
simple and $\nu$-separated, the bound is $C_\nu\varepsilon$.
No separation between roots of different components is required.
\end{lemma}
\begin{proof}
For $\varepsilon<\alpha_*/2$, the leading coefficient of $D_f+sE$
is invertible for every $0\le s\le1$. For sufficiently large $|z|$,
$|f_b(z)|\ge c|z|^d$ and $\|E(z)\|\le C\varepsilon|z|^d$.
A Neumann series excludes zeros outside a fixed disk uniformly in $s$.
Inside this disk, with $\rho=\dist(z,\Lambda)$,
\[
 \|D_f(z)^{-1}\|\le\alpha_*^{-1}\rho^{-d},\qquad
 \|E(z)\|\le C\varepsilon.
\]
Thus every zero throughout the homotopy lies in the union of disks of
radius $r=(C_1\varepsilon)^{1/d}$ about the true roots. Choose $C_1$
large enough that exclusion is strict outside these disks.

Fix $r'>r$ avoiding the finitely many radii at which components of
the disk union first touch. Distinct components of the closed $r'$-disk
union have positive distance from each other. Surround each by
piecewise smooth contours in the complement of the closed $r$-disks,
enclosing that component and no other. Such contours follow from a
slight outer perturbation of the disk-union boundary; holes may be
filled because all zeros are already in the $r$-disk union. The
homotopy determinant is nonzero on each contour. The argument principle
therefore preserves the enclosed number of roots with algebraic
multiplicity, equal to the number of true centres in that component.

If this number is $n$, a chain of overlapping disks joins centres in
at most $n-1$ steps. Their mutual distance is at most $2(n-1)r'$.
Every perturbed root is within $r$ of a centre, and there are exactly
$n$ such roots. Any bijection within the component has distance at most
$(2n-1)r'$. Let $r'\downarrow r$ through nonexceptional radii. Since
there are finitely many permutations, a subsequence has one fixed
matching. This proves the claim even at touching disks and repeated
centres. For real spectra, uncrossing inverted pairs gives ordered
matching.

Under the within-component separation assumption, the other factors
near a root of $f_b$ have magnitude bounded below in terms of $\nu$.
Away from these neighbourhoods $|f_b|$ is uniformly positive. On the
fixed disk,
$|f_b(z)|\ge c_\nu\min\{1,\dist(z,Z(f_b))\}$.
The same argument uses radius $C_\nu\varepsilon$. A repeated aggregate
root from different simple scalar components is a semisimple collision;
no eigenvector labeling or cross-component separation is needed.
\end{proof}

\begin{proof}[Proof of Theorem~\ref{thm:main}]
The defect, fibre and worst-case clock assertions have been proved.
The full-span assertion is proved in Proposition~\ref{prop:full-span}
below. A competitor satisfies $\cN_{P'}a'=b_{P'}$. Projection gives
\begin{equation}\label{eq:normalization-difference}
 \cN_P(a'-a)=(\Omega\otimes I)
       \big(q'(T_j)\operatorname{vec}(P(T_j)-P'(T_j))\big)_j.
\end{equation}
Uniform clock-value bounds for $q'$ yield $\|a'-a\|\le C\delta/\tau$;
compactness gives the cap in \eqref{eq:q-bound}.

Let $L_U=U^\dagger$ and $R_V=(V^\dagger)^{\mathsf T}$, so
$\|L_U\|\|R_V\|=\kappa^{-1}$. Set $M=L_UK'R_V$. At the first
$d+1$ clocks,
\begin{align}
 M(T_j)-\diag(\alpha_bf_b(T_j))
 &=q'(T_j)L_U(P'(T_j)-P(T_j))R_V\nonumber\\
 &\quad+(q'(T_j)-q(T_j))\diag(p_b(T_j)).
                                                        \label{eq:diagonal-error}
\end{align}
The second term is diagonal. Interpolation gives
$\|M-D_f\|_{\rm coef}\le C\delta(\kappa^{-1}+\tau^{-1})=:C\epsilon$.
In particular the normalization error has not acquired another channel
factor. For small $\epsilon$ the leading coefficient of $M$ is
invertible. It equals
$(L_UU')\diag(\alpha')(V'^{\mathsf T}R_V)$, so both square outer
factors are invertible. Hence
\[
 \det M(z)=\det(L_UU')\det(V'^{\mathsf T}R_V)
                           \prod_b\alpha_b'f_b'(z).
\]
Its roots are precisely the competitor multiset; no initial rank
hypothesis on that competitor was used. Lemma~\ref{lem:roots} gives
both small-error bounds. Bounded target range handles larger errors.
\end{proof}

\begin{remark}
Common-factor cancellation and multiplicity within a scalar component
are different mechanisms. The exponent $1/d$ is a uniform bound, not
a claim of the optimal exponent at every real-rooted singularity.
Theorem~\ref{thm:multiplicity} supplies matching neighbourhood rates
on explicitly stated multiplicity strata. Cross-component collisions
of simple scalar roots remain in the linear regime.
\end{remark}
